\documentclass[11pt]{amsart}

\usepackage[T1]{fontenc}
\usepackage{lmodern}
\usepackage{microtype}
\usepackage{amsmath,amssymb,amsthm}
\usepackage[hidelinks]{hyperref}

\hypersetup{
  pdftitle={Canon permutations avoiding 213 and 312 are counted by four Catalan numbers}
}

\newtheorem{theorem}{Theorem}
\newtheorem{lemma}[theorem]{Lemma}
\newtheorem{corollary}[theorem]{Corollary}
\theoremstyle{definition}
\newtheorem{remark}[theorem]{Remark}
\newtheorem*{conjecture*}{Conjecture (Laudone)}

\newcommand{\Cat}{C}
\newcommand{\wrd}[1]{\texttt{#1}}

\title[Canon permutations avoiding $213$ and $312$]
{Canon permutations avoiding $213$ and $312$ are counted by four Catalan
numbers}
\date{}

\subjclass[2020]{Primary 05A05; Secondary 05A15, 05A19, 05E10}
\keywords{canon permutation, pattern avoidance, Catalan number, ballot
sequence, standard Young tableau, multiset permutation}

\begin{document}

\begin{abstract}
Let $C_3^k$ be the set of canon permutations of the multiset
$\{1^k,2^k,3^k\}$, and let $c_3^k(\Lambda)$ count those avoiding every
pattern in $\Lambda$. Laudone conjectured that
$c_3^k(213,312)=4\Cat_k$ for all $k\ge 1$, where $\Cat_k$ is the $k$-th
Catalan number. We prove it, in the refined form
$c_3^{k,\sigma}(213,312)=\Cat_k$ for each of the four canon
witnesses $123$, $132$, $231$, $321$, and $0$ for the other two, by
exhibiting the surviving permutations explicitly: each is a solid block
$1^k$ at one end of the word and a ballot word on the letters $2$ and $3$ at
the other; the two choices of end and the two choices of which of $2,3$
dominates the ballot condition are what make four families of $\Cat_k$
words and not two. The proof is a six-case analysis of the canon witness and uses
the hypothesis $n=3$ in exactly one place, namely that over a three-letter
alphabet an occurrence of a three-letter pattern is a literal subsequence.
The same observation shows that relabelling the alphabet is a bijection of
$C_3^k$ transporting avoidance, whence $c_3^k(123,321)$ and
$c_3^k(132,231)$ equal $4\Cat_k$ as well, for every $k$; and that
$c_3^k(\{\sigma\})$ does not depend on which $\sigma\in S_3$ is forbidden.
The cells $k=1$ and $k=2$ of the conjecture are already in print.
\end{abstract}

\maketitle

\section{The conjecture}

For $n,k\ge1$ let $S_n^k$ be the set of permutations of the multiset
$\{1^k,2^k,\dots,n^k\}$: the words of length $nk$ that contain each letter
of $[n]$ exactly $k$ times. For a letter $a$ and $j\in[k]$ write
$\mathrm{pos}_a(j)$ for the position of the $j$-th occurrence of $a$. A word
$\pi\in S_n^k$ is \emph{canon} if there is a permutation $\sigma\in S_n$
such that, for every $j\in[k]$, the subsequence of $j$-th copies of the
letters read left to right spells $\sigma$; equivalently
\begin{equation}\label{eq:canon}
\mathrm{pos}_{\sigma_1}(j)<\mathrm{pos}_{\sigma_2}(j)<\cdots
<\mathrm{pos}_{\sigma_n}(j)
\quad\text{for all }j\in[k].
\end{equation}
Reading \eqref{eq:canon} at $j=1$ shows that $\sigma$ is the order in which the
letters first occur, so $\sigma$ is determined by $\pi$; we call it the
\emph{witness} of $\pi$ and write $C_n^k$ for the set of canon permutations and
$C_n^{k,\sigma}$ for those with witness $\sigma$. As usual $\pi$
\emph{contains} the pattern $p\in S_m$ if some length-$m$ subsequence of $\pi$
has the same relative order as $p$, and \emph{avoids} $p$ otherwise; for a set
$\Lambda\subseteq S_m$ we put
\begin{align*}
c_n^k(\Lambda)&=\#\{\pi\in C_n^k:\pi\text{ avoids every }p\in\Lambda\},\\
c_n^{k,\sigma}(\Lambda)&=\#\{\pi\in C_n^{k,\sigma}:
   \pi\text{ avoids every }p\in\Lambda\}.
\end{align*}
We write $\Cat_k=\frac{1}{k+1}\binom{2k}{k}$, so
$\Cat_1,\Cat_2,\Cat_3,\dots=1,2,5,14,42,132,\dots$; this is the indexing the
source itself uses.

Laudone \cite{Laudone} asks (his Question~6.5) for $c_3^k(\Lambda)$ as
$\Lambda$ ranges over subsets of $S_3$, with $n=3$ fixed and $k$ varying, and
offers one instance as an example; it is the last mathematical statement of
that paper, and we quote it verbatim. It appears in the final section, and is
numbered Conjecture~6.6 in the compiled source.

\begin{conjecture*}
For $k\ge1$, $c_3^k(213,312)=4C_k$ where $C_k$ is the $k$-th Catalan number.
\end{conjecture*}

Our results are the following. Theorem~\ref{thm:main} proves the conjecture
and refines it by the witness; Theorem~\ref{thm:orbit} shows that the same
argument settles two further pattern pairs, and Corollary~\ref{cor:single} one
further family.

\begin{theorem}\label{thm:main}
For every $k\ge1$,
\[
c_3^{k,\sigma}(213,312)=
\begin{cases}
\Cat_k,&\sigma\in\{123,132,231,321\},\\
0,&\sigma\in\{213,312\},
\end{cases}
\]
and consequently $c_3^k(213,312)=4\Cat_k$.
Moreover the survivors are exactly the following four families of words, where
$u$ always denotes a word on $k$ letters $2$ and $k$ letters $3$:
\begin{enumerate}
\item[\rm(i)] $1^k u$ with the $j$-th $2$ of $u$ before its $j$-th $3$ for all
      $j$ \quad (these have witness $123$);
\item[\rm(ii)] $1^k u$ with the $j$-th $3$ before the $j$-th $2$ \quad
      (witness $132$);
\item[\rm(iii)] $u\,1^k$ with the $j$-th $2$ before the $j$-th $3$ \quad
      (witness $231$);
\item[\rm(iv)] $u\,1^k$ with the $j$-th $3$ before the $j$-th $2$ \quad
      (witness $321$).
\end{enumerate}
\end{theorem}

\begin{theorem}\label{thm:orbit}
For every $\tau\in S_3$ and every $\Lambda\subseteq S_3$ we have
$c_3^k(\Lambda)=c_3^k(\tau\Lambda)$, where
$\tau\Lambda=\{\tau\circ p:p\in\Lambda\}$. The orbit of $\{213,312\}$ under
this action is
\[
\bigl\{\ \{213,312\},\ \{123,321\},\ \{132,231\}\ \bigr\},
\]
so for every $k\ge1$
\[
c_3^k(213,312)=c_3^k(123,321)=c_3^k(132,231)=4\Cat_k .
\]
\end{theorem}

\begin{corollary}\label{cor:single}
$c_3^k(\{\sigma\})$ is the same for all six $\sigma\in S_3$.
\end{corollary}

The cells $k=1$ and $k=2$ of the conjecture are already in print. At $k=1$ it
says that four of the six permutations of $S_3$ avoid $213$ and $312$. At
$k=2$ the canon permutations are precisely the nonnesting permutations of
Elizalde and Luo \cite{EL}, as \cite{Laudone} itself remarks, and their theorem
$c_n^2(132,231)=2^n$ gives $c_3^2(132,231)=8$; the complementation
$x\mapsto4-x$ is the relabelling $\tau=321$ of Theorem~\ref{thm:orbit}, which
carries $\{213,312\}$ to $\{132,231\}$, so $c_3^2(213,312)=8=4\Cat_2$ follows
from \cite{EL}. The argument below is elementary and short; we make no claim of
difficulty.

\section{Containment over three letters is literal}

Everything below rests on one observation, and it is the only place where the
hypothesis $n=3$ enters.

\begin{lemma}\label{lem:literal}
Let $w$ be a word over the alphabet $\{1,2,3\}$ and let $p\in S_3$. Then $w$
contains $p$ if and only if $p_1p_2p_3$ occurs in $w$ as a literal subsequence,
that is, if and only if there are positions $i<j<l$ with $w_i=p_1$, $w_j=p_2$,
$w_l=p_3$.
\end{lemma}

\begin{proof}
An occurrence of $p$ is a triple of positions $i<j<l$ whose values
$(a,b,c)=(w_i,w_j,w_l)$ have the same relative order as $p$. Since $p$ has
three distinct entries, $a,b,c$ are pairwise distinct; being three distinct
elements of $\{1,2,3\}$ they are $1,2,3$ in some order, and having the same
relative order as $p$ then forces $(a,b,c)=(p_1,p_2,p_3)$. The converse is
immediate.
\end{proof}

In particular, over $\{1,2,3\}$: $w$ contains $213$ iff some $2$ precedes some
$1$ which precedes some $3$, and $w$ contains $312$ iff some $3$ precedes some
$1$ which precedes some $2$. We use these two sentences without further
comment.

\begin{remark}\label{rem:n4}
Lemma~\ref{lem:literal} is false for larger alphabets, and $\wrd{314}$ is a
witness: it contains $213$ (the values $3,1,4$ have the relative order of
$2,1,3$) but has no literal $2,1,3$. This is exactly why nothing in this note
says anything about $c_n^k(213,312)$ for $n\ge4$.
\end{remark}

\section{Proof of Theorem~\ref{thm:main}}

Fix $k\ge1$ and write $\Lambda=\{213,312\}$. Since the witness is determined by
the word, the six sets $C_3^{k,\sigma}$ partition $C_3^k$, and it suffices to
treat each witness.

\smallskip\noindent
\textbf{(a) $\sigma=213$ and $\sigma=312$ are empty.}
If $\sigma=213$ then \eqref{eq:canon} at $j=1$ gives
$\mathrm{pos}_2(1)<\mathrm{pos}_1(1)<\mathrm{pos}_3(1)$, so the letters
$2,1,3$ occur in that order and $\pi$ contains $213$ by
Lemma~\ref{lem:literal}. If $\sigma=312$ then
$\mathrm{pos}_3(1)<\mathrm{pos}_1(1)<\mathrm{pos}_2(1)$ and $\pi$ contains
$312$. Hence $c_3^{k,213}(\Lambda)=c_3^{k,312}(\Lambda)=0$.

\smallskip\noindent
\textbf{(b) $\sigma=123$.}
Let $\pi\in C_3^{k,123}$ avoid $\Lambda$. We claim all $k$ letters $1$ come
first. If not, some letter of $\{2,3\}$ precedes some $1$; let $q$ be the
leftmost position of a $1$ having a non-$1$ to its left, say
$q=\mathrm{pos}_1(j)$.
\begin{itemize}
\item If a $2$ occurs before $q$: by \eqref{eq:canon},
      $\mathrm{pos}_3(j)>\mathrm{pos}_2(j)>\mathrm{pos}_1(j)=q$, so some $3$
      occurs after $q$; that $2$, the $1$ at $q$ and that $3$ give an
      occurrence of $213$, a contradiction.
\item If no $2$ but some $3$ occurs before $q$: then
      $\mathrm{pos}_3(1)<q<\mathrm{pos}_2(1)$, contradicting
      $\mathrm{pos}_2(1)<\mathrm{pos}_3(1)$, which is \eqref{eq:canon} for
      $\sigma=123$.
\end{itemize}
So $\pi=1^ku$ with $u$ a word on $k$ letters $2$ and $k$ letters $3$. For such
a word $\mathrm{pos}_1(j)=j$ is smaller than every position of $u$, so
\eqref{eq:canon} reduces to $\mathrm{pos}_2(j)<\mathrm{pos}_3(j)$, i.e.\ the
$j$-th $2$ of $u$ precedes its $j$-th $3$, for every $j$. Conversely every such
word avoids $\Lambda$: an occurrence of $213$ or of $312$ needs a $1$ strictly
to the right of a $2$ or of a $3$, and there is none. Reading $2$ as an up-step
and $3$ as a down-step identifies these $u$ with the Dyck paths of semilength
$k$, since ``the $j$-th $2$ precedes the $j$-th $3$ for all $j$'' says exactly
that every prefix of $u$ has at least as many $2$'s as $3$'s; the Catalan count
of those paths is standard, see \cite{Stanley}. Hence
$c_3^{k,123}(\Lambda)=\Cat_k$, realised by family~(i).

\smallskip\noindent
\textbf{(c) $\sigma=231$.}
Let $\pi\in C_3^{k,231}$ avoid $\Lambda$ and let $q=\mathrm{pos}_1(1)$. By
\eqref{eq:canon}, $\mathrm{pos}_2(1)<\mathrm{pos}_3(1)<q$, so both a $2$ and a
$3$ occur before $q$. If some $3$ occurred after $q$ we would have
$2\cdots1\cdots3$ and hence the pattern $213$; if some $2$ occurred after $q$
we would have $3\cdots1\cdots2$ and hence $312$. So every position after $q$
carries a $1$, i.e.\ $\pi=u\,1^k$ with $u$ a word on $k$ letters $2$ and $k$
letters $3$. Now $\mathrm{pos}_1(j)=2k+j$ exceeds every position of $u$, so
\eqref{eq:canon} again reduces to ``the $j$-th $2$ of $u$ precedes its $j$-th
$3$ for every $j$'', and conversely each such word avoids $\Lambda$ because no
$2$ or $3$ follows a $1$. Hence $c_3^{k,231}(\Lambda)=\Cat_k$, realised by
family~(iii).

\smallskip\noindent
\textbf{(d) $\sigma=132$ and $\sigma=321$.}
Let $t$ be the relabelling that exchanges the letters $2$ and $3$ and fixes
$1$. Applying $t$ letterwise is an involution of $S_3^k$; by \eqref{eq:canon}
it carries $C_3^{k,\sigma}$ onto $C_3^{k,t\circ\sigma}$, and by
Lemma~\ref{lem:literal} it carries occurrences of $213$ to occurrences of $312$
and back. Since $t$ therefore permutes $\Lambda$, it restricts to a bijection
between the $\Lambda$-avoiders with witness $\sigma$ and those with witness
$t\circ\sigma$. As $t\circ 123=132$ and $t\circ 231=321$, cases (b) and (c)
give $c_3^{k,132}(\Lambda)=c_3^{k,321}(\Lambda)=\Cat_k$, and applying $t$ to
families (i) and (iii) yields (ii) and (iv).

\smallskip\noindent
Summing, $c_3^k(\Lambda)=0+0+4\Cat_k=4\Cat_k$. \qed

\section{Proof of Theorem~\ref{thm:orbit} and Corollary~\ref{cor:single}}

Let $\tau\in S_3$ act on a word letterwise. This is a bijection of $S_3^k$; by
\eqref{eq:canon} a canon word with witness $\sigma$ is sent to a canon word
with witness $\tau\circ\sigma$, so $\tau$ is a bijection of $C_3^k$. By
Lemma~\ref{lem:literal}, containment of $p\in S_3$ is containment of the
literal word $p_1p_2p_3$, and $\tau$ carries a literal occurrence of
$p_1p_2p_3$ to a literal occurrence of $\tau(p_1)\tau(p_2)\tau(p_3)$ and back.
Hence $\pi$ avoids $\Lambda$ if and only if $\tau(\pi)$ avoids $\tau\Lambda$,
and $c_3^k(\Lambda)=c_3^k(\tau\Lambda)$.

Writing $\tau$ as the word $\tau(1)\tau(2)\tau(3)$, one computes
$\tau\{213,312\}$ for the six $\tau$ and finds $\{213,312\}$ for
$\tau\in\{123,132\}$, $\{123,321\}$ for $\tau\in\{213,231\}$ and
$\{132,231\}$ for $\tau\in\{312,321\}$. So the stabiliser is the order-two
group generated by the transposition of $2$ and $3$ and the orbit has size
three, as stated; Theorem~\ref{thm:main} then gives the common value $4\Cat_k$.

For Corollary~\ref{cor:single}, $S_3$ acts simply transitively on the
singletons $\{\sigma\}$, since $\tau\circ p=p'$ has the unique solution
$\tau=p'p^{-1}$. Hence all six values $c_3^k(\{\sigma\})$ coincide. \qed

\section{Verification}

A companion program \texttt{verify.py} (Python~3.9+, standard library only,
exact integer arithmetic) re-derives the quantities asserted here. It builds
$C_3^k$ twice --- once from definition~\eqref{eq:canon} over all of $S_3^k$ and
once from standard Young tableaux --- decides containment twice, once by
relative order and once literally, and confirms the refinement of
Theorem~\ref{thm:main} and the four families by exhaustive enumeration for
$k\le6$; a transfer-matrix recursion that uses only the prefix form of
\eqref{eq:canon} and a greedy pattern matcher, and knows nothing of
Theorem~\ref{thm:main}, reproduces the refinement and the value $4\Cat_k$ for
all $k\le24$ and, for the witness $123$ alone, for all $k\le40$. It also
verifies Lemma~\ref{lem:literal} exhaustively over the alphabet, exhibits its
failure at $n=4$, and checks Theorem~\ref{thm:orbit} object by object for
$k\le3$. No $k$ beyond those ranges, no statement about $n\ge4$ and no
bibliographic assertion is machine-checked; the proof above is what carries the
general case.

\begin{thebibliography}{9}

\bibitem{EL}
S.~Elizalde and A.~Luo,
\emph{Pattern avoidance in nonnesting permutations},
Discrete Math. Theor. Comput. Sci. \textbf{27} (2025), no.~1
(Permutation Patterns 2024 special issue), published 17 October 2025;
\href{https://doi.org/10.46298/dmtcs.14885}{doi:10.46298/dmtcs.14885};
\href{https://arxiv.org/abs/2412.00336}{arXiv:2412.00336}.

\bibitem{Laudone}
R.~Laudone,
\emph{Pattern avoidance in canon permutations},
preprint, 21 August 2026;
\href{https://arxiv.org/abs/2608.21351}{arXiv:2608.21351}.

\bibitem{Stanley}
R.~P. Stanley,
\emph{Catalan Numbers},
Cambridge University Press, Cambridge, 2015.

\end{thebibliography}

\end{document}
